\phantomsection\addcontentsline{toc}{section}{Marine Renewable Energy}
\ECchapter{25}{Marine Renewable Energy}{Power from oceans}{ECBlue}

\ECObjectives{
\begin{itemize}
\item Compare tidal, wave and offshore energy technologies.
\item Explain the mechanical and environmental design constraints.
\item Estimate power, annual energy yield and capacity factor.
\end{itemize}}

\begin{multicols}{2}
\begin{engineeringnews}
Oceans contain major renewable energy resources. Offshore wind is already deployed at commercial scale,
whereas tidal-stream and wave-energy devices are still developing. Marine systems benefit from strong,
partly predictable resources but face demanding installation, grid-connection and maintenance conditions.
\end{engineeringnews}

\begin{ecarticle}
Marine renewable energy includes offshore wind, tidal barrages, tidal-stream turbines and wave-energy
converters. Each technology extracts energy from a different physical phenomenon and therefore requires a
different structure, control strategy and maintenance plan.

Tidal-stream turbines operate like underwater wind turbines. Because water is roughly 800 times denser
than air, even a moderate current can exert large forces. The available kinetic power can be estimated by
\[
P_{\mathrm{avail}}=\frac{1}{2}\rho A v^3,
\]
where $\rho$ is water density, $A$ the swept area and $v$ the current speed. The useful electrical power is
lower because the rotor, generator and converter have limited efficiencies. Tides are highly predictable,
but suitable sites are limited and the equipment must withstand strong cyclic loads.

Wave-energy converters use the relative motion between waves and a floating or fixed structure. Some
devices use oscillating water columns, while others use hinged bodies or hydraulic power take-off systems.
Wave direction, period and height vary continuously, so the controller must adapt while preventing excessive
motion. Survival during an extreme storm is often more demanding than normal energy production.

The marine environment accelerates corrosion and biological growth. Salt water attacks materials, while
barnacles and algae modify surface properties and increase drag. Seals, cables and bearings must operate for
long periods with limited access. Retrieving a device for repair may require a specialised vessel and a calm
weather window, so reliability strongly affects the levelised cost of energy.

Engineers must also consider navigation, fishing, underwater noise and marine ecosystems. Environmental
monitoring is required before and after installation. Subsea cables and offshore substations can make grid
connection expensive. The best design is therefore not simply the device with the highest peak power, but
the system that delivers dependable energy over its complete life cycle.
\end{ecarticle}

\begin{tcolorbox}[title=\sffamily\bfseries Did You Know?,colback=yellow!10,colframe=ECOrange]
Because tidal motion is driven by astronomical cycles, its timing can be predicted years in advance. Tidal
energy is variable, but far more predictable than wind or solar generation.
\end{tcolorbox}

\begin{engineeringfocus}
\textbf{Tidal-stream conversion chain}
\begin{center}
\resizebox{0.96\linewidth}{!}{%
\begin{tikzpicture}[font=\scriptsize,>=Latex,node distance=3mm]
\node[draw=ECBlue,fill=ECLightBlue,rounded corners] (flow) {tidal current};
\node[draw=ECGreen,fill=ECGreen!10,rounded corners,right=of flow] (rotor) {rotor};
\node[draw=ECPurple,fill=ECPurple!10,rounded corners,right=of rotor] (gen) {generator};
\node[draw=ECOrange,fill=ECOrange!10,rounded corners,right=of gen] (conv) {converter};
\node[draw=ECRed,fill=ECRed!8,rounded corners,below=of gen] (grid) {subsea cable and grid};
\draw[->,thick] (flow)--(rotor); \draw[->,thick] (rotor)--(gen); \draw[->,thick] (gen)--(conv);
\draw[->,thick] (conv)--(grid);
\end{tikzpicture}}
\end{center}
\end{engineeringfocus}

\begin{keyfigures}
\begin{tabularx}{\linewidth}{@{}Xr@{}}
Predictable resource & tides\\
Main material threat & corrosion\\
Common surface issue & biofouling\\
Major cost driver & offshore maintenance\\
Grid connection & subsea cable
\end{tabularx}
\end{keyfigures}

\begin{tcolorbox}[title=\sffamily\bfseries Illustrative capacity factors,colback=white,colframe=ECBlue]
\begin{center}
\begin{tikzpicture}
\begin{axis}[xbar,width=0.96\linewidth,height=48mm,xlabel={Capacity factor (\%)},xmin=0,xmax=55,
symbolic y coords={Tidal barrage,Offshore wind,Wave energy,Tidal stream},ytick=data,
grid=major,ticklabel style={font=\scriptsize},nodes near coords]
\addplot table[x=factor,y=technology,col sep=comma]{data/EC25/capacity.csv};
\end{axis}
\end{tikzpicture}
\end{center}
\footnotesize Values are illustrative; real performance depends strongly on the site and technology maturity.
\end{tcolorbox}

\begin{tcolorbox}[title=\sffamily\bfseries Quantitative application,colback=ECBlue!4,colframe=ECBlue]
A tidal turbine has a rotor diameter of $18\,\mathrm{m}$ in a current of $2.5\,\mathrm{m\,s^{-1}}$.
Take $\rho=1025\,\mathrm{kg\,m^{-3}}$, a power coefficient $C_p=0.40$ and an electrical efficiency of
$90\%$.
\begin{enumerate}
\item Calculate the swept area.
\item Estimate the electrical power produced at this current speed.
\item If the device operates at an average capacity factor of $35\%$, estimate its annual energy output.
\end{enumerate}
\end{tcolorbox}

\begin{vocabulary}
\begin{tabularx}{\linewidth}{@{}lX@{}}
tidal stream & courant de marée\\
wave period & période de vague\\
power take-off & système de conversion d'énergie\\
cyclic load & charge cyclique\\
corrosion & corrosion\\
biofouling & encrassement biologique\\
subsea cable & câble sous-marin\\
capacity factor & facteur de charge\\
mooring & amarrage\\
storm survival & tenue en tempête
\end{tabularx}
\end{vocabulary}

\begin{questions}
\item Why can relatively slow seawater flows produce large forces?
\item What makes wave-energy control difficult?
\item Why does maintenance strongly influence marine energy costs?
\item Give two environmental impacts that must be monitored.
\item Why is peak power insufficient for comparing devices?
\end{questions}

\begin{engineeringchallenge}
Choose a tidal-stream site and propose a 1 MW device. Estimate the required rotor area for an assumed
current speed, select corrosion protection, define a maintenance strategy and identify two ecological
monitoring measures.
\end{engineeringchallenge}

\begin{speaking}
Should marine renewable projects be prioritised in protected coastal areas when they reduce carbon
emissions? Discuss biodiversity, energy security, local jobs and precaution.
\end{speaking}
\end{multicols}

\subsection*{Sources}
\ECsource{International Energy Agency -- Ocean Power}{https://www.iea.org/energy-system/renewables/ocean-power}
\ECsource{European Marine Energy Centre}{https://www.emec.org.uk/}